% infini_vs_fini_0.tex
% Formalisation fidèle de l'hypothèse philosophique « infini vs fini »
% Version 0 — document de travail
\documentclass[11pt,a4paper]{article}

\usepackage[utf8]{inputenc}
\usepackage[T1]{fontenc}
\usepackage[french]{babel}
\usepackage{lmodern}
\usepackage{microtype}
\usepackage{amsmath,amssymb,amsthm,mathtools}
\usepackage{geometry}
\usepackage{enumitem}
\usepackage{booktabs}
\usepackage{longtable}
\usepackage{hyperref}
\usepackage{csquotes}

\geometry{margin=2.6cm}
\setlist{itemsep=0.35em, topsep=0.5em}
\hypersetup{colorlinks=true, linkcolor=blue!45!black, urlcolor=blue!45!black}

\newtheorem{axiom}{Axiome}
\newtheorem{definition}{Définition}
\newtheorem{proposition}{Proposition}
\newtheorem{remark}{Remarque}
\newtheorem{question}{Question ouverte}

\newcommand{\Inf}{\mathcal I}
\newcommand{\Fin}{\mathcal F}
\newcommand{\Lim}{\mathsf{Lim}}
\newcommand{\Det}{\mathsf{Det}}
\newcommand{\Cons}{\mathsf{Cons}}

\title{\textbf{Infini vs fini}\\
\large Formalisation fidèle d'une hypothèse métaphysique}
\author{Document de travail --- version 0}
\date{23 septembre 2026}

\begin{document}
\maketitle

\begin{abstract}
Ce document présente et organise l'hypothèse selon laquelle l'infini est fondamental, d'abord au sens ontologique, tandis que ce que nous appelons ordinairement « fini » apparaît comme une détermination produite par des contraintes. L'objectif est la fidélité conceptuelle : les catégories standard du fini et de l'infini ne sont pas prises comme prémisses neutres permettant de juger l'hypothèse ; elles sont précisément parmi les notions que l'hypothèse remet en question.

Le document distingue strictement trois niveaux : (i) la thèse métaphysique elle-même ; (ii) les définitions humaines, nécessairement limitées, par lesquelles nous tentons de la penser ; (iii) les formalismes mathématiques ou scientifiques usuels, qui peuvent être comparés à la thèse mais ne doivent pas la remplacer. La présente version ne prétend pas démontrer la vérité de la thèse. Elle vise à la rendre aussi explicite, rigoureuse et non-trahie que possible, tout en signalant les points où un argument supplémentaire serait nécessaire.
\end{abstract}

\tableofcontents

\section{Statut et méthode du document}

\subsection{Objet}

L'hypothèse défendue dans ce document commence par un refus d'une certaine symétrie intuitive entre les notions de « fini » et d'« infini ». Dans l'usage courant, on peut les traiter comme deux propriétés déjà définies : certaines choses seraient finies, d'autres infinies. La thèse propose au contraire de remonter d'un niveau et de demander si cette opposition est elle-même conceptuellement légitime.

Le point de départ est donc l'inversion suivante : le fini n'est pas posé comme réalité première, puis comparé à l'infini ; c'est la non-limitation qui est prise comme fondement, et les formes finies comme résultats de limitations ou de déterminations.

\subsection{Principe de non-réduction}

Le principe méthodologique central est :

\begin{quote}
\emph{Ne pas juger l'hypothèse depuis une définition standard du fini et de l'infini lorsque cette définition est précisément l'objet de la contestation.}
\end{quote}

Cela ne signifie pas que les mathématiques, la physique ou les autres usages standards seraient invalides. Cela signifie seulement qu'ils ne doivent pas être importés clandestinement comme axiomes métaphysiques de départ.

\subsection{Objectivité interne}

L'objectivité recherchée ici consiste à distinguer ce qui découle de l'hypothèse, ce qui est encore une intuition, ce qui constitue une conséquence conditionnelle, et ce qui demeure une question ouverte. Une proposition ne sera pas déclarée démontrée simplement parce qu'elle est cohérente avec le cadre.

\section{Thèse ontologique fondamentale}

La formulation centrale est :

\begin{quote}
\textbf{L'infini est fondamental ; le fini émerge comme résultat de contraintes, de délimitations ou de déterminations.}
\end{quote}

Le mot « fondamental » est ici prioritairement ontologique. Il ne signifie pas seulement que le concept d'infini serait utile pour faire des mathématiques, ni que l'idée d'infini serait psychologiquement antérieure. Il signifie que, dans l'ordre de réalité postulé par la thèse, la non-limitation est première et la finitude dérivée.

Le schéma ontologique proposé est donc :

\[
\boxed{
\Inf
\;\longrightarrow\;
\text{contrainte}
\;\longrightarrow\;
\text{détermination}
\;\longrightarrow\;
\Fin
}
\]

La flèche ne désigne pas nécessairement une succession temporelle. Elle désigne une dépendance ontologique proposée : le fini n'est pas premier ; il résulte d'une limitation ou d'une détermination.

\section{Redéfinition du « fini »}

\subsection{Refus du fini absolu}

La thèse ne nie pas qu'une description puisse comporter un nombre déterminé d'éléments ou qu'une opération puisse s'arrêter. Ce qu'elle conteste est le passage de ce constat local à l'affirmation d'une finitude absolue.

Ainsi, écrire

\[
1,2,3
\]

est compatible avec le fait que l'on se soit arrêté à \(3\). Mais la proposition :

\[
\text{« 3 est la fin en soi »}
\]

introduit une affirmation beaucoup plus forte. La question « fin de quoi ? » devient alors fondamentale. Une fin ne semble identifiable qu'en relation avec un domaine dans lequel une frontière est posée.

\subsection{Exemple 1 : \(1,2,3\)}

Dans le cadre de la thèse, \(123\) ne doit pas être défini ontologiquement comme « quelque chose qui s'arrête à 3 ». Sa définition pertinente est :

\[
\boxed{123 = \text{une détermination particulière au sein de l'infini}.}
\]

Le symbole \(123\) n'est pas ici déclaré inexistant ou dénué de propriétés. Il représente une forme déterminée, choisie, délimitée ou distinguée. Ce qui est refusé est le statut de « fin ontologique ultime » attribué à cette détermination.

On peut représenter intuitivement :

\[
\underbrace{1,2,3}_{\text{détermination}}
\;\big|
\;4,5,6,\ldots
\]

où la barre exprime une coupure opérée par le cadre de considération.

Le caractère important de l'exemple est que l'arrêt à \(3\) n'est pas, à lui seul, une démonstration d'une frontière intrinsèque de l'être.

\subsection{Formulation générale}

La notion de fini est alors mieux exprimée par :

\[
\boxed{\Fin_x = \Det_x(\Inf)}
\]

où \(x\) représente le cadre ou l'ensemble de contraintes considéré et \(\Det_x\) une opération de détermination. Cette notation est une formalisation conceptuelle et non une théorie mathématique achevée.

La proposition essentielle est donc :

\[
\boxed{\text{fini} = \text{détermination sous contrainte},}
\]

et non :

\[
\boxed{\text{fini} = \text{réalité possédant intrinsèquement une limite absolue}.}
\]

\section{Redéfinition de « l'infini »}

\subsection{Définition propre à l'hypothèse}

L'infini ne doit pas être compris comme une grandeur extrêmement grande. Il est plutôt défini par sa non-limitabilité.

Formulation provisoire :

\[
\boxed{
\Inf := \text{ce qui ne peut être limité, totalisé ni définitivement contraint}.
}
\]

Cette définition est intentionnellement négative. Elle ne prétend pas produire une image ou une représentation exhaustive de l'infini ; elle indique ce que l'on refuse de lui attribuer : une borne ultime, une clôture définitive, un extérieur qui constituerait sa limite ou une contrainte capable de l'enfermer complètement.

\subsection{Pas une « chose infiniment grande »}

Dans ce cadre, l'infini n'est pas nécessairement un objet dont la grandeur serait simplement supérieure à celle de tous les objets finis. La notion de grandeur n'est pas primitive dans la thèse. Ce qui est fondamental est l'impossibilité d'imposer une limite ultime.

Il faut donc éviter l'équivalence :

\[
\text{infini} = \text{très grand}.
\]

La thèse soutient plutôt :

\[
\text{infini} = \text{non-limitable au sens ontologique proposé}.
\]

\subsection{L'infini est inconcevable au sens exhaustif}

Un point central de l'hypothèse est la distinction entre l'infini et notre concept de l'infini.

Une conscience, un langage et des moyens matériels eux-mêmes limités ne peuvent pas, selon la thèse, produire une représentation absolument exhaustive d'une réalité qui ne peut elle-même être limitée ou totalisée.

On doit donc distinguer :

\[
\boxed{
\text{infini tel qu'il est supposé être} \neq \text{concept humain de l'infini}.
}
\]

Cela n'implique pas que toute pensée de l'infini soit impossible. Cela signifie que toute définition humaine de l'infini est une détermination partielle produite depuis une condition limitée.

\subsection{La définition de l'infini est elle-même limitée}

La thèse comprend une forme de réflexivité : notre propre notion de l'infini est déjà une contrainte.

En effet, nous parlons de « l'infini », nous le distinguons, nous lui associons des propriétés et nous lui donnons un mot. Toute cette opération appartient à notre système cognitif et linguistique.

Ainsi :

\[
\boxed{
\text{notion humaine de l'infini} = \text{détermination humaine limitée de ce que nous appelons l'infini}.
}
\]

La définition proposée n'est donc pas présentée comme l'infini lui-même. Elle est une approximation conceptuelle rendue possible par nos propres limites.

\section{Pourquoi « totalité » est problématique}

Une objection classique consiste à demander : « Si tout est une partie de l'infini, qu'est-ce que la totalité de l'infini ? »

Dans le cadre de cette hypothèse, la question est mal formée si « totalité » signifie déjà une clôture, une complétude ou une frontière finale.

La raison est structurelle :

\[
\text{totaliser} \;\Rightarrow\; \text{déterminer une limite de ce qui est totalisé},
\]

alors que :

\[
\Inf \;\Rightarrow\; \text{non-limitabilité et non-totalisabilité}.
\]

Dire « totalité de l'infini » revient donc à appliquer à l'infini une opération que la définition même de l'hypothèse exclut.

\boxed{
\text{Il n'existe pas de « totalité » de l'infini si « totalité » signifie une clôture ultime.}
}

Cette conclusion ne signifie pas que le concept de totalité est dépourvu de sens dans tous les contextes. Il signifie seulement qu'il ne peut pas être introduit sans précaution comme caractéristique d'une réalité définie comme non-totalisable.

\section{Pas d'extérieur absolu à l'infini}

Si l'on prétend que quelque chose \(X\) est extérieur à \(\Inf\), on pose une distinction :

\[
\Inf \quad | \quad X.
\]

Cette distinction constitue une frontière conceptuelle. Or la définition de l'infini interdit une limite ultime.

On obtient donc la conséquence conditionnelle :

\[
\boxed{
\text{il n'existe pas d'extérieur absolu de }\Inf.
}
\]

Cette proposition ne veut pas dire que toutes les choses particulières sont identiques. Elle affirme que la différence entre deux déterminations particulières ne doit pas être transformée automatiquement en une séparation absolue d'avec l'infini.

\section{Le statut de « partie de l'infini »}

L'expression « partie de l'infini » doit être comprise avec une prudence particulière. Elle ne signifie pas nécessairement que l'infini soit un ensemble au sens mathématique standard contenant des sous-ensembles.

Son rôle dans l'hypothèse est métaphysique :

\[
\boxed{
\text{toute détermination finie est une détermination interne au domaine non-limitable}.
}
\]

Ainsi, « \(123\) est une partie de l'infini » signifie d'abord :

\begin{quote}
la détermination « 123 » ne possède pas de statut ontologique de clôture absolue ; elle est une forme localement délimitée, exprimée depuis l'intérieur d'une réalité que l'hypothèse considère comme non-limitable.
\end{quote}

Il ne faut donc pas réintroduire clandestinement les propriétés d'un ensemble totalisant pour formaliser cette relation.

\section{L'exemple des dix-sept particules}

Considérons un cadre hypothétique dans lequel un univers est décrit comme contenant exactement dix-sept particules :

\[
U = \{p_1,\dots,p_{17}\}.
\]

La réponse standard peut être : « \(U\) est fini. »

Dans le cadre de l'hypothèse, cette réponse n'est pas suffisante pour établir une finitude ontologique absolue. Elle établit seulement qu'un certain cadre a été défini avec dix-sept éléments.

Le nombre dix-sept constitue alors une détermination de ce cadre. La question métaphysique n'est pas :

\[
17 \text{ est-il le nombre d'éléments de } U\ ?
\]

mais :

\[
\text{Que signifie ontologiquement le fait que } U \text{ soit arrêté à }17 ?
\]

Si l'on dit qu'une dix-huitième particule n'existe pas, l'hypothèse ne répond pas automatiquement qu'elle existe. Elle répond que l'absence de cette dix-huitième particule dans ce cadre ne suffit pas à établir que « 17 » soit une frontière absolue de l'être.

Autrement dit :

\[
\boxed{
\text{« exactement 17 dans un cadre » } \neq \text{ « 17 est la limite absolue de la réalité ».}
}
\]

\section{Le caractère central des contraintes}

\subsection{Définition provisoire}

Une contrainte est ici tout ce qui transforme une possibilité non ultimement bornée en une détermination particulière.

Le mot « contrainte » ne doit pas être réduit aux contraintes physiques. Il peut désigner, selon le contexte :

\begin{itemize}
    \item une restriction matérielle ;
    \item une condition physique ;
    \item une capacité cognitive ;
    \item une convention de description ;
    \item une frontière spatiale ou temporelle ;
    \item une règle de sélection ;
    \item une distinction conceptuelle.
\end{itemize}

La thèse générale est que ces mécanismes produisent ou manifestent des formes finies.

\subsection{Attention : toutes les contraintes ne sont pas nécessairement du même type}

Pour rester rigoureux, il faut séparer :

\[
\text{contrainte épistémique} \neq \text{contrainte physique} \neq \text{contrainte ontologique}.
\]

L'hypothèse affirme que la finitude telle qu'elle se présente dans les déterminations dépend de contraintes, mais elle n'a pas encore démontré que toutes les contraintes observées appartiennent à une même catégorie fondamentale.

Cette distinction est essentielle pour éviter de confondre les limites de notre connaissance avec les limites de la réalité.

\section{Différence entre ignorance et impossibilité ontologique}

L'énoncé :

\[
\text{« nous ne pouvons pas dépasser }X\text{ »}
\]

est différent de :

\[
\text{« il est ontologiquement impossible de dépasser }X\text{ »}.
\]

La thèse vise la seconde forme lorsque celle-ci est utilisée pour caractériser une frontière absolue.

Exemple : le fait qu'un être humain ne puisse pas compter toutes les étoiles observables ne signifie pas que leur nombre soit infini. Une limite de connaissance n'est pas automatiquement une non-limitation ontologique.

Cette différence ne détruit pas la thèse ; elle précise son niveau de discours.

\section{Le problème de la « première contrainte »}

La thèse rencontre un problème structurel important : si le fini provient des contraintes, quelle est l'origine de la première contrainte ?

Trois structures sont concevables sans qu'aucune soit encore démontrée.

\subsection{Régression}

Une contrainte peut dépendre d'une autre :

\[
C_1 \leftarrow C_2 \leftarrow C_3 \leftarrow \cdots
\]

Cette possibilité maintient une chaîne sans première contrainte. Elle serait compatible avec la structure générale de la thèse, mais elle demande une théorie précise de cette dépendance.

\subsection{Première contrainte}

Une première contrainte \(C_0\) pourrait être postulée :

\[
\Inf \longrightarrow C_0 \longrightarrow \Fin.
\]

Mais alors il faut expliquer pourquoi \(C_0\) existe, pourquoi cette contrainte particulière est fondamentale et en quel sens elle peut limiter une réalité définie comme non-limitable.

\subsection{Absence de première contrainte}

Il est également possible que l'idée même d'une « première » contrainte soit issue de notre manière finie de penser en termes de commencement, d'origine et de point initial.

Cette possibilité est cohérente avec la thèse mais n'en constitue pas encore une conséquence démontrée.

\section{Statut du réel, du possible et du conceptualisé}

L'hypothèse accorde une place importante au fait que la distinction entre le réel et le conceptualisé est elle-même formulée par un sujet situé et limité.

Cela ne permet toutefois pas de conclure automatiquement :

\[
\text{réel} = \text{imaginaire}.
\]

La thèse plus fidèle est :

\[
\boxed{
\text{la distinction réel / imaginé est une distinction produite dans un cadre ; elle ne doit pas être confondue avec une frontière ontologique absolue.}
}
\]

Le statut exact du possible, de l'imaginé, de l'actuel et du matériel reste à développer.

\section{Pourquoi la conception humaine de l'infini est nécessairement imparfaite}

L'hypothèse ne dit pas : « nous ne pouvons rien dire de l'infini ». Elle dit plutôt : « ce que nous pouvons en dire est nécessairement une détermination humaine ».

On peut distinguer :

\[
\text{concevoir une propriété de l'infini}
\]

et

\[
\text{contenir l'infini dans une représentation complète}.
\]

La seconde opération est incompatible avec la définition proposée si « contenir complètement » signifie totaliser ou limiter.

Ainsi, une définition comme

\[
\Inf := \text{ce qui ne peut être ultimement limité}
\]

est à comprendre comme une approximation conceptuelle négative : elle dit ce que l'on refuse d'attribuer à l'infini, sans prétendre en fournir une image exhaustive.

\section{Rapport aux mathématiques}

\subsection{Ce qui doit être évité}

Il serait incorrect de partir d'une définition mathématique du type :

\[
\{1,2,3\} \text{ est fini}
\]

puis de conclure immédiatement que l'hypothèse métaphysique est réfutée. Une telle démarche suppose précisément que la définition mathématique du « fini » possède une portée ontologique absolue.

\subsection{Ce qui peut être admis}

Les mathématiques peuvent définir rigoureusement des objets finis et infinis dans leurs propres cadres formels. Cela ne constitue pas, en soi, une thèse sur la constitution ultime du réel.

La distinction pertinente est donc :

\[
\boxed{
\text{validité interne d'un formalisme} \neq \text{démonstration ontologique sur la réalité}.
}
\]

La thèse ne demande pas de nier les mathématiques usuelles ; elle refuse qu'une formalisation particulière soit utilisée sans argument comme définition ultime de l'être.

\subsection{Le symbole \texorpdfstring{$\infty$}{infty}}

Le symbole \(\infty\) est une représentation humaine. Il ne doit pas être identifié sans reste avec l'infini au sens ontologique de l'hypothèse :

\[
\boxed{\infty_{\mathrm{symbole}} \neq \Inf_{\mathrm{ontologique}}.}
\]

Le premier appartient à un langage ; le second désigne ce que le langage tente, nécessairement imparfaitement, de pointer.

\section{Critique de la fausse symétrie « fini / infini »}

La structure standard suivante est suspecte dans le cadre de la thèse :

\[
\Fin \quad \leftrightarrow \quad \Inf.
\]

Elle suggère deux catégories ontologiques de même rang.

La structure proposée est au contraire asymétrique :

\[
\boxed{
\Inf \longrightarrow \Det(\Inf) = \Fin_x.
}
\]

Le fini dépend d'une détermination ; l'infini, lui, n'est pas défini comme une chose qui dépend du fini pour exister.

Cela constitue l'un des éléments d'originalité de l'hypothèse : le fini n'est plus l'opposé de l'infini ; il est une forme limitée issue d'une réalité non-limitable.

\section{Le principe d'inconcevabilité sans contradiction performative}

Une objection possible serait : « si l'infini est inconcevable, comment peux-tu en parler ? »

La réponse de la thèse est qu'il faut distinguer l'objet visé et la représentation utilisée pour le viser.

Une expression peut désigner une limite de ce qui peut être représenté sans constituer une représentation exhaustive de cette limite.

Dans ce sens :

\[
\boxed{
\text{parler de l'infini} \neq \text{posséder une représentation totale de l'infini}.
}
\]

La thèse n'est donc pas autocontradictoire simplement parce qu'elle utilise le concept « infini » pour affirmer que l'infini dépasse toute conceptualisation exhaustive.

\section{Structure logique provisoire de l'hypothèse}

Le noyau de la théorie peut être présenté par les propositions suivantes.

\begin{axiom}[Primauté ontologique]
L'infini est fondamental dans l'ordre ontologique proposé.
\end{axiom}

\begin{axiom}[Non-limitabilité]
Aucune limite ontologique ultime ne constitue l'infini.
\end{axiom}

\begin{axiom}[Non-totalisabilité]
L'infini ne peut pas être transformé en totalité close sans introduire une notion étrangère à sa définition.
\end{axiom}

\begin{axiom}[Finitude dérivée]
Ce qui est appelé « fini » est une détermination produite sous contrainte, et non une finitude absolue présupposée comme fondamentale.
\end{axiom}

\begin{axiom}[Non-identité entre représentation et réalité]
Aucune représentation humaine de l'infini ne doit être identifiée sans reste à l'infini lui-même.
\end{axiom}

\begin{axiom}[Absence d'extérieur absolu]
L'infini ne peut pas être traité comme une région limitée par un extérieur ontologique absolu.
\end{axiom}

Ces axiomes ne constituent pas encore une démonstration ; ils constituent la base déclarée de l'hypothèse.

\section{Conséquences conditionnelles}

Si les axiomes précédents sont admis, plusieurs conséquences suivent au niveau conceptuel.

\begin{proposition}[Non-fondamentalité du fini]
Le fini ne possède pas, dans ce cadre, le statut de catégorie ontologique première.
\end{proposition}

\begin{proof}[Justification conditionnelle]
Par l'axiome de primauté ontologique, l'infini est fondamental. Par l'axiome de finitude dérivée, les formes finies résultent de déterminations sous contrainte. Leur statut est donc dérivé plutôt que premier.
\end{proof}

\begin{proposition}[Non-absoluité des limites locales]
Le fait qu'une structure soit bornée dans un cadre donné ne suffit pas, à lui seul, à établir l'existence d'une limite ontologique absolue.
\end{proposition}

\begin{proof}[Justification conditionnelle]
Une borne décrite dans un cadre est une détermination relative à ce cadre. Pour devenir une limite absolue, il faudrait établir qu'aucune détermination plus large ou aucune distinction d'un autre ordre ne peut être formulée. Cette étape supplémentaire n'est pas fournie par la simple constatation de la borne locale.
\end{proof}

\begin{proposition}[Réflexivité de la conception humaine]
La représentation humaine de l'infini est elle-même une détermination finie.
\end{proposition}

\begin{proof}[Justification conditionnelle]
La représentation est produite par des capacités cognitives, linguistiques et matérielles particulières. Elle est donc située dans un cadre déterminé et ne constitue pas, selon l'axiome de non-identité, l'infini lui-même.
\end{proof}

\section{Ce que l'hypothèse n'affirme pas}

Pour éviter toute trahison ou tout élargissement non justifié, la présente version ne prétend pas avoir établi les propositions suivantes :

\begin{enumerate}
    \item que l'espace physique observable ou l'Univers physique est effectivement infini ;
    \item que le temps est nécessairement éternel dans le passé ou dans le futur ;
    \item qu'il existe effectivement un nombre infini de particules ou d'objets physiques ;
    \item que toute possibilité imaginée est physiquement réelle ;
    \item que toute distinction entre réel et imaginaire est ontologiquement nulle ;
    \item que les mathématiques standard sont incohérentes ;
    \item que le concept mathématique standard de fini est faux dans son propre domaine formel ;
    \item que l'absence actuelle d'une possibilité implique sa réalisation future ;
    \item qu'il existe nécessairement une « matière infinie » ou une « énergie infinie » ;
    \item que l'hypothèse est déjà démontrée par les exemples de \(123\) ou de dix-sept particules.
\end{enumerate}

Ces points restent soit ouverts, soit hors du périmètre de la thèse telle qu'elle a été formulée.

\section{Objections sérieuses à traiter ultérieurement}

La fidélité à l'hypothèse n'exige pas de masquer ses difficultés. Les objections suivantes sont pertinentes sans qu'elles permettent à elles seules de la déclarer fausse.

\subsection{Objection 1 : possibilité d'une réalité absolument finie}

On peut demander s'il est logiquement contradictoire qu'une réalité soit absolument bornée. La thèse doit expliquer pourquoi une limite absolue serait impossible, et pas seulement difficile à imaginer.

\subsection{Objection 2 : glissement entre possible et réel}

Du fait qu'une extension puisse être conçue ou décrite, il ne s'ensuit pas automatiquement qu'elle soit réellement possible ou réelle. La théorie doit préciser le rapport entre possibilité, actualité et conceptualisation.

\subsection{Objection 3 : ambiguïté du mot « contrainte »}

Si toute finitude provient d'une contrainte, il faut déterminer si la contrainte est :

\[
\text{cognitive},\quad \text{descriptive},\quad \text{physique},\quad \text{logique},\quad \text{ontologique},
\]

ou si une catégorie plus fondamentale doit être introduite.

\subsection{Objection 4 : origine des contraintes}

Si une première contrainte existe, elle demande une explication. Si aucune première contrainte n'existe, la théorie doit expliquer la structure de la régression ou de l'absence de premier terme.

\subsection{Objection 5 : auto-référence}

Si toute définition de l'infini est elle-même une limitation, la définition de l'infini doit être traitée comme une approximation et non comme une essence capturée. La théorie semble devoir accepter cette limitation sans considérer cela comme une réfutation.

\section{Distinctions indispensables pour la suite}

Pour que la théorie reste cohérente, plusieurs couples conceptuels doivent être maintenus séparés :

\[
\begin{array}{c@{\qquad}c}
\text{limite locale} & \neq \text{limite absolue}\\
\text{ignorance} & \neq \text{non-limitabilité ontologique}\\
\text{définition} & \neq \text{réalité définie}\\
\text{représentation} & \neq \text{objet visé}\\
\text{possibilité conceptuelle} & \neq \text{existence réelle}\\
\text{distinction} & \neq \text{séparation ontologique absolue}.
\end{array}
\]

Ces distinctions sont importantes non pour ramener la théorie aux standards, mais pour empêcher que sa propre logique ne glisse d'une proposition à une autre sans justification.

\section{Formulation synthétique de la thèse}

La formulation la plus compacte et la plus fidèle à ce stade est :

\begin{quote}
\textbf{L'infini est la réalité supposée fondamentalement non-limitable, non-totalisable et sans extérieur absolu. Le « fini » n'est pas une propriété ontologique première ; il désigne des déterminations locales issues de contraintes ou de découpages. Toute représentation humaine de l'infini est elle-même une détermination limitée et ne peut donc prétendre être l'infini dans sa totalité.}
\end{quote}

On peut également écrire le mouvement conceptuel :

\[
\boxed{
\Inf
\rightarrow
\text{contrainte}
\rightarrow
\text{frontière}
\rightarrow
\text{distinction}
\rightarrow
\text{détermination}
\rightarrow
\Fin
}
\]

mais sans réifier cette notation : l'infini, dans l'hypothèse, n'est précisément pas une « chose » au sens ordinaire qu'une définition exhaustive pourrait enfermer.

\section{Point de départ pour une version ultérieure}

La prochaine étape théorique devrait examiner la question :

\[
\boxed{
\text{Si l'infini est fondamental, comment et pourquoi apparaissent les contraintes ?}
}
\]

Cette question est centrale parce que l'hypothèse ne veut pas simplement affirmer deux faits indépendants : « l'infini existe » et « il existe des choses finies ». Elle veut expliquer le second à partir du premier.

Une théorie plus développée devra donc clarifier :

\begin{enumerate}
    \item ce qu'est exactement une contrainte ;
    \item comment une contrainte produit une détermination ;
    \item si toute détermination implique nécessairement une frontière ;
    \item si plusieurs types de contraintes doivent être distingués ;
    \item si une première contrainte est nécessaire ou non ;
    \item comment la pluralité et la différence émergent sans introduire une limite absolue ;
    \item dans quel sens cette métaphysique pourrait être mise en rapport avec les mathématiques et les sciences sans être absorbée par leurs catégories propres.
\end{enumerate}

\section*{Conclusion}

L'hypothèse étudiée n'est pas une simple affirmation quantitative selon laquelle « il y aurait énormément de choses ». Elle est une proposition sur le statut ontologique des limites elles-mêmes.

Son intuition directrice est :

\[
\boxed{
\text{ce que nous appelons « fini » est une détermination ; ce qui est fondamental est le non-limitable.}
}
\]

La définition de l'infini demeure volontairement non-exhaustive : puisque toute conception humaine est elle-même limitée, l'infini ne peut être contenu sans reste dans le concept que nous en formons. Cette limitation de notre propre définition ne constitue donc pas une anomalie de la théorie ; elle fait partie de sa structure.

La rigueur demandée par cette hypothèse consiste dès lors moins à imposer les définitions standards qu'à éviter les glissements : ne pas transformer une borne locale en borne absolue, une définition en réalité, une ignorance en infinité, une possibilité conceptuelle en existence et une totalité en concept neutre.

\bigskip
\begin{center}
\textit{Version 0 : formalisation fidèle, ouverte à la critique et à la consolidation argumentative.}
\end{center}

\end{document}
